\chapter{Matériels et méthodes}
\label{chap:methodes}

Ce chapitre présente le cadre épistémologique, les variables, l'environnement technique et le
protocole expérimental retenus pour évaluer la fonction \texttt{computeRFM} d'ELMES-TEKA.
La présentation suit le cadre \emph{ADEMP} proposé par \citet{Morris2019}, dont les cinq
composantes — Objectifs, Mécanismes générateurs de données, Cibles, Méthodes et Mesures de
performance — structurent l'ensemble du dispositif. Le détail opérationnel est consigné à
l'Annexe~\ref{ann:protocole}.

% ------------------------------------------------------------------
\section{Positionnement épistémologique : Design Science Research, posture pragmatique
et évaluation par simulation}
\label{sec:methodes-epistemologie}

\subsection{Design Science Research comme stratégie globale}
\label{subsec:methodes-dsr}

La recherche s'inscrit dans le paradigme du \emph{Design Science Research}
(DSR) tel que formalisé par % TODO: \citep{Hevner2004}.
Dans ce cadre, la connaissance scientifique est produite par la construction et
l'évaluation d'artefacts destinés à résoudre des problèmes pratiques. ELMES-TEKA
constitue l'artefact central de la démarche : conçu pour répondre aux besoins de gestion
commerciale des petits entrepreneurs congolais, il intègre un moteur de segmentation RFM
dont la validité algorithmique doit être établie avant tout usage en conditions réelles.

Le DSR se déploie en cinq étapes successives : (1)~la conscience du problème, soit
l'identification des risques associés à un calcul de segmentation erroné dans un contexte
multi-tenant ; (2)~la suggestion, matérialisée par la fonction \texttt{computeRFM} en
TypeScript ; (3)~le développement de l'artefact, soit l'instrumentation du banc de
simulation ; (4)~l'\textbf{évaluation}, qui constitue la contribution principale du présent
mémoire ; et (5)~la conclusion, apportant les implications pour la pratique et la recherche.
Le mémoire se concentre donc sur l'étape évaluative du cycle DSR.

\subsection{Posture pragmatique}
\label{subsec:methodes-pragmatisme}

La posture épistémologique adoptée est pragmatique : la connaissance produite vaut par
son utilité et par la capacité de l'artefact à résoudre les problèmes concrets des
commerçants. Il ne s'agit pas d'établir une vérité universelle sur les algorithmes RFM,
mais de déterminer si \texttt{computeRFM} se comporte de manière fiable dans les
conditions d'usage prévues par ELMES-TEKA. Cette posture justifie le recours à des
corpus synthétiques construits pour couvrir des configurations représentatives et des
cas limites, plutôt qu'à des données de terrain dont la collecte dépasserait le périmètre
du mémoire.

\subsection{Raisonnement hypothético-déductif}
\label{subsec:methodes-hypotheses}

La vérification des deux hypothèses de recherche procède par raisonnement
hypothético-déductif, en cohérence avec une posture post-positiviste. Deux propositions
falsifiables sont formulées et soumises à des tests empiriques sur données synthétiques :

\begin{itemize}
  \item \textbf{H1 (conformité)} : la fonction TypeScript \texttt{computeRFM} produit des
    résultats identiques à l'implémentation Python de référence, pour tout corpus relevant
    du domaine couvert par la campagne.
  \item \textbf{H2 (robustesse)} : les clients dont les valeurs RFM sont éloignées des
    seuils de terciles ne changent pas de segment lorsqu'une perturbation monétaire de
    $\pm1\,\%$ est appliquée à leurs montants de commandes.
\end{itemize}

Chaque hypothèse est opérationnalisée par un critère de décision quantitatif dont le seuil
est fixé avant l'exécution de la campagne, conformément au principe de préenregistrement
préconisé par \citet{Lakens2013}.

\subsection{Cadre ADEMP de présentation de la simulation}
\label{subsec:methodes-ademp}

La simulation est présentée selon le cadre \emph{ADEMP} proposé par \citet{Morris2019} et
détaillé dans \citet{Morris2019slides}. Ce cadre décompose toute étude de simulation en cinq
composantes explicites :

\begin{enumerate}
  \item \textbf{Aims (Objectifs)} : déterminer si \texttt{computeRFM} est conforme à la
    référence indépendante (H1) et si la segmentation est robuste à une perturbation
    monétaire modérée pour les clients éloignés des seuils (H2).
  \item \textbf{Data-generating mechanisms (Mécanismes générateurs)} : trois profils de
    génération de corpus — régulier (distribution uniforme discrète), dissymétrique
    (distributions exponentielle et log-normale), identique (valeurs constantes) — croisés
    avec quatre tailles de corpus ($n \in \{3, 10, 50, 200\}$ clients).
  \item \textbf{Estimands (Cibles)} : taux de non-conformité inter-implémentations (H1) et
    taux de changement de segment parmi les clients éloignés sous bruit monétaire $\pm1\,\%$
    (H2).
  \item \textbf{Methods (Méthodes)} : comparaison exacte au centime USD entre le moteur
    JavaScript \texttt{computeRFM} et l'oracle Python indépendant ; classification des
    clients éloignés des seuils selon la règle à $5\,\%$ décrite au
    paragraphe~\ref{subsec:methodes-critere-h2}.
  \item \textbf{Performance measures (Mesures de performance)} : taux de non-conformité
    global, taux de changement de segment sur clients éloignés, borne supérieure de
    l'intervalle de confiance Monte-Carlo à $95\,\%$, taux de clients perdus
    (\texttt{lost\_rate}) et indice de Rand ajusté (\texttt{ari}).
\end{enumerate}

\subsection{Limites épistémiques}
\label{subsec:methodes-limites-epistemiques}

Les corpus utilisés dans cette étude sont synthétiques et ne sont pas calibrés sur des
données issues d'entreprises congolaises réelles. Les résultats décrivent le comportement
de \texttt{computeRFM} sur les mécanismes générateurs retenus ; ils ne permettent pas de
conclure directement sur le comportement attendu en production ni sur les distributions
réelles des données des commerçants congolais. Cette limite constitue une frontière
épistémique inhérente à toute évaluation par simulation préalable au déploiement terrain.

% ------------------------------------------------------------------
\section{Variables d'étude}
\label{sec:methodes-variables}

Le tableau~\ref{tab:ann-dictionnaire} (Annexe~\ref{ann:protocole}) présente le dictionnaire
complet des variables. Les éléments essentiels sont résumés ci-dessous.

\paragraph{Variables indépendantes (VI).}
Les variables indépendantes définissent les conditions expérimentales :

\begin{itemize}
  \item \textbf{Taille du corpus} : $n \in \{3, 10, 50, 200\}$ clients par essai ; cette
    variation permet d'observer l'effet de la masse de données sur la stabilité des terciles.
  \item \textbf{Profil de génération} : trois modalités — \texttt{regular} (uniforme discret
    sur les montants et les fréquences), \texttt{skewed} (exponentielle pour les fréquences,
    log-normale pour les montants), \texttt{ties} (valeurs identiques pour tous les clients,
    induisant des ex~æquo systématiques dans les calculs de terciles).
  \item \textbf{Type de perturbation} : neuf variantes appliquées à chaque corpus initial —
    bruit monétaire ($\pm1\,\%$, $\pm5\,\%$, $\pm20\,\%$), décalage de dates ($\pm1$~jour,
    $\pm7$~jours), injection de commandes dupliquées ($1\,\%$, $5\,\%$ du nombre de
    commandes), suppression de commandes ($1\,\%$, $5\,\%$ du nombre de commandes).
\end{itemize}

\paragraph{Variables dépendantes (VD).}
\begin{itemize}
  \item \textbf{Taux de non-conformité} (H1) : proportion d'essais dans lesquels au moins
    un client présente un écart entre le moteur JavaScript et l'oracle Python, que ce soit
    sur l'un des trois scores ($s_R$, $s_F$, $s_M$) ou sur l'étiquette de segment.
  \item \textbf{Taux de changement de segment sur clients éloignés} (H2) : proportion de
    clients éloignés des seuils de terciles dont le segment diffère entre la partition initiale
    et la partition après perturbation.
\end{itemize}

\paragraph{Variables de contrôle (VC).}
\begin{itemize}
  \item Graine maître : \texttt{20261102} (date de référence de la campagne).
  \item Générateur pseudo-aléatoire : Mulberry32, déterministe et portable entre
    environnements.
  \item Date d'observation fixe : 1\textsuperscript{er}~octobre~2026 à 00:00:00~UTC.
  \item Devise : USD exclusivement.
\end{itemize}

La table de configuration complète figure au tableau~\ref{tab:ann-configuration}
(Annexe~\ref{ann:protocole}).

% ------------------------------------------------------------------
\section{Environnement technique et versions}
\label{sec:methodes-environnement}

\begin{table}[htb]
  \centering
  \caption{Environnement technique de la campagne de simulation.}
  \label{tab:methodes-environnement}
  \begin{tabular}{ll}
    \toprule
    Composant & Version \\
    \midrule
    Node.js        & 24.22.1 \\
    TypeScript     & 5.9.3 \\
    Python         & 3.12.14 \\
    NumPy          & 2.3.5 \\
    pandas         & 2.2.3 \\
    SciPy          & 1.17.0 \\
    Matplotlib     & 3.10.8 \\
    Commit applicatif & \texttt{\small b062a93a7be96d082dee9afe66b9f4f1f9ebecaa} \\
    \bottomrule
  \end{tabular}
\end{table}

La fonction évaluée est \texttt{computeRFM}, telle qu'elle existe dans le dépôt
ELMES-TEKA au commit \texttt{b062a93a7be96d082dee9afe66b9f4f1f9ebecaa}. Aucune
modification fonctionnelle n'est apportée au code applicatif pendant la campagne ;
seule une couche d'instrumentation externe est ajoutée pour intercepter les appels
et enregistrer les sorties.

L'oracle Python est une implémentation indépendante développée spécifiquement pour
cette étude. Il reproduit la même logique de calcul des terciles et les mêmes règles
d'attribution des segments RFM que la fonction TypeScript, mais en additionnant les
montants en centimes entiers afin d'éliminer les erreurs d'arrondi propres à
l'arithmétique à virgule flottante. L'oracle est développé sans partager de code avec
le moteur JavaScript ; son indépendance est une condition nécessaire à la validité de
la comparaison.

Les statuts de commandes retenus pour le calcul RFM sont : \texttt{CONFIRMED},
\texttt{PAID} et \texttt{DELIVERED}. Seules les commandes libellées en USD sont prises
en compte ; les commandes dans d'autres devises sont exclues du calcul sans lever
d'exception.

Les mécanismes de génération des données et les paramètres des distributions sont
décrits au tableau~\ref{tab:ann-generateurs} (Annexe~\ref{ann:protocole}).

% ------------------------------------------------------------------
\section{Protocole d'expérimentation}
\label{sec:methodes-protocole}

\subsection{Structure de la campagne}
\label{subsec:methodes-structure}

La campagne porte sur \textbf{12\,000~essais} résultant du croisement exhaustif des
conditions expérimentales suivantes :

\begin{itemize}
  \item 4 tailles de corpus : $n \in \{3, 10, 50, 200\}$ ;
  \item 3 profils de génération : \texttt{regular}, \texttt{skewed}, \texttt{ties} ;
  \item 100 répétitions par cellule, chacune tirée avec une graine distincte dérivée
    de la graine maître \texttt{20261102} ;
  \item 10 variantes par répétition : la situation initiale (sans perturbation) et
    9 perturbations.
\end{itemize}

La combinaison $4 \times 3 \times 100 \times 10 = 12\,000$ essais constitue le domaine
principal de la campagne. Cinq configurations supplémentaires hors domaine (devise mixte,
identifiant absent, date invalide, date future, montant négatif) ont été soumises à titre
exploratoire et sont traitées séparément au chapitre~\ref{chap:resultats}.

\subsection{Principe d'appariement des perturbations}
\label{subsec:methodes-appariement}

Chaque perturbation est appariée au corpus initial de la même répétition. Elle est
générée indépendamment à partir d'une graine de perturbation dérivée de la graine de
corpus, de sorte que deux perturbations distinctes appliquées au même corpus initial
ne se cumulent pas. Les seuils de terciles sont recalculés intégralement après chaque
perturbation ; ils ne sont jamais reportés d'une variante à l'autre. Ce principe garantit
que chaque essai constitue une comparaison entre deux états indépendants du même corpus
de départ.

\subsection{Mécanismes de génération des profils}
\label{subsec:methodes-generateurs}

\paragraph{Profil régulier (\texttt{regular}).}
Le nombre de commandes par client est tiré d'une distribution uniforme discrète sur
$[1;\, 20]$. Les montants unitaires sont tirés uniformément sur $[1{,}00;\, 100{,}00]$~USD,
arrondis au centime. Les dates de commande sont réparties uniformément dans les douze mois
précédant la date d'observation.

\paragraph{Profil dissymétrique (\texttt{skewed}).}
Le nombre de commandes suit une distribution exponentielle tronquée, produisant une
majorité de clients peu actifs et une minorité de clients très actifs. Les montants
suivent une distribution log-normale, générant des valeurs concentrées autour d'un faible
ticket moyen avec une queue épaisse vers les montants élevés. Ce profil vise à reproduire
les asymétries documentées dans les bases de données clients réelles
\citep{Fader2005, Salminen2023}.

\paragraph{Profil identique (\texttt{ties}).}
Tous les clients du corpus partagent exactement le même nombre de commandes (trois),
la même date de dernier achat (trente jours avant la date d'observation) et le même
montant de ticket (10,00~USD). Ce profil force des ex~æquo systématiques dans les trois
dimensions RFM, testant ainsi le comportement de la fonction dans les cas de dégénérescence
des terciles.

\paragraph{Commandes anciennes.}
Dans les profils régulier et dissymétrique, un sous-ensemble de commandes est généré
avec des dates antérieures à la fenêtre principale et des montants multipliés par un
facteur aléatoire compris dans $[0{,}80;\, 1{,}20]$, afin d'introduire de l'hétérogénéité
dans la dimension temporelle sans modifier la distribution globale des montants.

\subsection{Perturbations appliquées}
\label{subsec:methodes-perturbations}

Le tableau~\ref{tab:methodes-perturbations} récapitule les neuf perturbations appliquées
à chaque corpus initial.

\begin{table}[htb]
  \centering
  \caption{Les neuf perturbations appliquées à chaque corpus initial.}
  \label{tab:methodes-perturbations}
  \small
  \begin{tabular}{llp{0.45\textwidth}}
    \toprule
    Identifiant & Dimension & Description \\
    \midrule
    \texttt{amount\_noise\_1pct}  & Monétaire  & Multiplication de chaque montant par un facteur uniforme sur $[0{,}99;\, 1{,}01]$ \\
    \texttt{amount\_noise\_5pct}  & Monétaire  & Multiplication par un facteur uniforme sur $[0{,}95;\, 1{,}05]$ \\
    \texttt{amount\_noise\_20pct} & Monétaire  & Multiplication par un facteur uniforme sur $[0{,}80;\, 1{,}20]$ \\
    \texttt{date\_noise\_1d}      & Temporelle & Décalage aléatoire de $\pm1$~jour sur chaque date de commande \\
    \texttt{date\_noise\_7d}      & Temporelle & Décalage aléatoire de $\pm7$~jours sur chaque date de commande \\
    \texttt{duplicates\_1pct}     & Structure  & Duplication aléatoire de $1\,\%$ des commandes existantes \\
    \texttt{duplicates\_5pct}     & Structure  & Duplication aléatoire de $5\,\%$ des commandes existantes \\
    \texttt{missing\_1pct}        & Structure  & Suppression aléatoire de $1\,\%$ des commandes existantes \\
    \texttt{missing\_5pct}        & Structure  & Suppression aléatoire de $5\,\%$ des commandes existantes \\
    \bottomrule
  \end{tabular}
\end{table}

% ------------------------------------------------------------------
\section{Critères de conformité et de robustesse}
\label{sec:methodes-criteres}

\subsection{Critère de conformité H1}
\label{subsec:methodes-critere-h1}

Un essai est déclaré \textbf{non conforme} dès qu'au moins un client présente, pour cet
essai, un écart entre le moteur JavaScript et l'oracle Python sur l'un des éléments
suivants :

\begin{itemize}
  \item le score de récence $s_R \in \{1, 2, 3\}$ ;
  \item le score de fréquence $s_F \in \{1, 2, 3\}$ ;
  \item le score de montant $s_M \in \{1, 2, 3\}$ ;
  \item l'étiquette de segment parmi les cinq catégories : «~Champions~»,
    «~Loyaux~», «~À risque~», «~Petits acheteurs~», «~Potentiels~».
\end{itemize}

La comparaison est exacte au centime USD : deux sommes identiques en représentation
décimale à deux chiffres sont considérées égales, quelle que soit leur représentation
interne en virgule flottante. L'oracle Python réalise toutes les sommes de montants en
centimes entiers avant de diviser par 100 pour obtenir la valeur en USD ; cette convention
constitue la référence.

Le taux de non-conformité est défini par :
\[
  \tau_{\mathrm{H1}} = \frac{\text{nombre d'essais non conformes}}{\text{nombre total d'essais}}.
\]

La décision de compatibilité avec H1 n'est pas fondée sur un seuil arbitraire de
$\tau_{\mathrm{H1}}$ ; il s'agit au contraire de caractériser la nature et la source
des écarts observés, conformément à l'objectif diagnostique du mémoire.

\subsection{Critère de robustesse H2}
\label{subsec:methodes-critere-h2}

\paragraph{Définition des clients éloignés.}
Soient $q_{d,j}$ le $j$-ème quantile de la dimension $d \in \{R, F, M\}$, avec $j \in
\{1, 2\}$ correspondant aux deux seuils de tercile (tercile inférieur et tercile
supérieur). Un client $i$ est dit \textbf{éloigné} de tous les seuils si, pour chacune
des six paires $(d, j)$, la distance normalisée suivante est strictement supérieure à
$5\,\%$ :

\begin{equation}
  \frac{\bigl|x_{i,d} - q_{d,j}\bigr|}{\max\!\bigl(|q_{d,j}|,\, 1\bigr)} > 0{,}05,
  \quad \forall\, d \in \{R, F, M\},\; j \in \{1, 2\}.
  \label{eq:critere-eloignement}
\end{equation}

Cette règle garantit que les clients retenus pour évaluer H2 ne sont pas proches d'une
frontière de segmentation, de sorte qu'une perturbation mineure ne devrait pas les faire
changer de tercile — et donc de segment — sous l'hypothèse de robustesse.

\paragraph{Cellules admissibles pour H2.}
Seules les cellules vérifiant les trois conditions suivantes sont admissibles pour le
critère H2 :

\begin{enumerate}
  \item taille de corpus $n \in \{50, 200\}$ ;
  \item profil \texttt{regular} ou \texttt{skewed} (le profil \texttt{ties} est exclu en
    raison de la dégénérescence de ses seuils de tercile) ;
  \item scénario de bruit monétaire $\pm1\,\%$ (\texttt{amount\_noise\_1pct}).
\end{enumerate}

De plus, une répétition n'est admissible que si les deux partitions (initiale et perturbée)
comportent au moins dix clients communs éloignés au sens de l'équation~\eqref{eq:critere-eloignement}. Une cellule est retenue si elle comprend au moins 80~répétitions
admissibles.

\paragraph{Indicateur et règle de décision.}
Le taux de changement de segment sur clients éloignés est calculé pour chaque répétition
admissible. Un intervalle de confiance Monte-Carlo à $95\,\%$ est construit par loi~$t$
à partir des répétitions admissibles d'une cellule, puis tronqué à $[0;\, 1]$. La
décision est déclarée \textbf{Compatible} avec H2 si la borne supérieure de cet intervalle
est strictement inférieure ou égale à $5\,\%$. Ce seuil est fixé avant l'exécution de la
campagne.

% ------------------------------------------------------------------
\section{Vérification préliminaire par cas manuels}
\label{sec:methodes-cas-manuels}

Avant le lancement de la campagne automatisée de 12\,000~essais, neuf cas manuels ont été
construits à la main et enregistrés dans le fichier \texttt{conformity.json}. Cette étape
préliminaire vise à valider l'oracle Python avant de l'utiliser comme référence pour
l'ensemble des essais automatisés.

Chaque cas manuel est conçu pour couvrir une configuration limite susceptible de
révéler un défaut de l'oracle ou une divergence de convention entre les deux
implémentations. Les neuf configurations couvertes sont les suivantes :

\begin{enumerate}
  \item corpus vide (aucun client, aucune commande) ;
  \item six clients aux frontières exactes des quantiles, couvrant les cinq groupes
    de segments en même temps ;
  \item client singleton (un seul client dans le corpus) ;
  \item valeurs toutes identiques (ex~æquo parfaits sur les trois dimensions) ;
  \item filtrage par statut de commande (\texttt{CONFIRMED}, \texttt{PAID},
    \texttt{DELIVERED}) — vérification que les commandes à d'autres statuts sont ignorées ;
  \item exclusion par tenant et exercice — vérification de l'isolation multi-tenant ;
  \item client présent dans le répertoire mais sans commande admissible ;
  \item montant nul sur l'ensemble des commandes d'un client ;
  \item date de commande avec fraction de journée — vérification du traitement de la
    partie entière pour le calcul de récence.
\end{enumerate}

Le tableau~\ref{tab:ann-cas-manuels} (Annexe~\ref{ann:protocole}) liste les entrées et
sorties attendues de chacun des neuf cas. La satisfaction de ces cas ne constitue pas une
preuve d'absence d'écart dans la campagne principale ; elle établit uniquement que l'oracle
Python et le moteur JavaScript s'accordent sur les configurations limites les plus
prévisibles, ce qui est une condition nécessaire à la validité de la comparaison
automatisée.
